% Generated draft: requires mathematical and source review.
\documentclass[11pt]{article}
\usepackage[margin=25mm]{geometry}
\usepackage{fontspec}
\setmainfont{Latin Modern Roman}
\usepackage{amsmath,amssymb,mathtools,unicode-math}
\setmathfont{Latin Modern Math}
\usepackage{enumitem,needspace,xurl,hyperref}
\setlength{\parindent}{0pt}
\setlength{\parskip}{5pt}
\setlength{\emergencystretch}{3em}
\newcommand{\N}{\mathbb{N}}
\newcommand{\Z}{\mathbb{Z}}
\newcommand{\Q}{\mathbb{Q}}
\newcommand{\R}{\mathbb{R}}
\newcommand{\C}{\mathbb{C}}
\newcommand{\F}{\mathbb{F}}
\newcommand{\A}{\mathbb{A}}
\newcommand{\PP}{\mathbb{P}}
\newcommand{\E}{\mathbb{E}}
\newcommand{\eps}{\varepsilon}
\newcommand{\dd}{\,\mathrm{d}}
\begin{document}
\section*{UnsolvedMath ID: 1101103}
\textbf{Source problem code:} catalogue entry
\subsection{Definitions and mathematical statement}
For metric spaces $X,Y$, a quasi-isometry is a map $f:X\to Y$ for which there exist $L\geq1$, $A\geq0$ such that
\[
L^{-1}d_X(x,x')-A\leq d_Y(f(x),f(x'))\leq Ld_X(x,x')+A
\quad(x,x'\in X),
\]
and every point of $Y$ is within distance $A$ of $f(X)$. Two maps $f,g:X\to Y$ represent the same coarse map if $\sup_{x\in X}d_Y(f(x),g(x))<\infty$. The group $\operatorname{QI}(X)$ consists of self-quasi-isometries modulo this relation, with composition as group operation.

Let $M_\alpha$ range over the Gromov--Thurston negatively pinched manifold examples intended by the source, and write $X_\alpha=\widetilde M_\alpha$ for their universal covers with lifted metrics. The record does not supply the construction, parameter range, dimensions, or precise metric hypotheses for this family. Under the usual universal-cover reading of the geometric-group-theory question, Kleiner's terse question has two related targets:
\begin{enumerate}
\item Determine $\operatorname{QI}(X_\alpha)$ for each example, including whether the coarse classes are forced by a specified geometric symmetry group.
\item Determine the equivalence relation on examples given by
\[
\alpha\sim\beta\quad\Longleftrightarrow\quad
\exists\text{ a quasi-isometry }X_\alpha\to X_\beta.
\]
\end{enumerate}
The source does not explicitly choose one of these readings or specify universal covers. The distinction is material: a question about the metric of a compact manifold itself has a different coarse meaning from the question about its universal cover or fundamental group. Neither the family nor the intended rigidity group can be reconstructed uniquely from the supplied sentence.

The supplied summary reports both self-quasi-isometry rigidity and quasi-isometry classification as open, citing Sisto--Viaggi (2026) for the latter's being poorly understood. It also reports its own proof-sketch of $X_k\not\simeq_{\mathrm{qi}}\mathbb H^d$ whenever $M_k$ admits no hyperbolic metric, invoking Tukia's uniform quasiconformal group theorem and Waldhausen/Farrell--Jones rigidity against the Gromov--Thurston examples. That sketch and its hypotheses are not reproduced in this record. Its further reported obstruction to a wall-pattern approach is the absence of quasi-wall rigidity and of an exponent-$1$ conclusion in Bourdon's M\"obius extension for snowflaked boundary maps. The summary says known distinctions among these manifolds are at the homotopy or commensurability level rather than the quasi-isometry level. These claims are retained as supplied context, not independently established conclusions.
\subsection{Short English statement}
Kleiner asks about quasi-isometries of the Gromov--Thurston examples. In the natural universal-cover setting this may ask for each self-quasi-isometry group, for which examples are quasi-isometric to one another, or both. The supplied summary reports both questions open, but the precise family and intended interpretation need the original source context.
\subsection{Sources}
\begin{itemize}
\item Local UnsolvedMath registry, permanent ID 1101103, exact source code catalogue entry, ``Questions in Geometric Group Theory --- Q 11.3''; status open, classification OPEN-TRIAGE.
\item Bestvina, \emph{Questions in Geometric Group Theory} (2004), Question 11.3, PDF page 19, attributed to Kleiner, \url{https://www.math.utah.edu/~bestvina/eprints/questions-updated.pdf}. The provenance describes manual transcription of a public author-hosted PDF, accessed July 29, 2026, and says its open status was not independently updated. The research summary qualifies its own verification: its full PDF fetch stopped before Section 11, and it confirmed the wording through a search-engine extract matching the worklist transcription.
\item Sisto--Viaggi (2026), cited by the supplied summary for the poorly understood quasi-isometry classification. No title, identifier, theorem number, or fuller reference is provided.
\item The summary invokes Gromov--Thurston, Tukia's uniform quasiconformal group theorem, Waldhausen/Farrell--Jones rigidity, and Bourdon's M\"obius extension in its reported analysis. No precise external citations or complete proof-sketch are supplied for these uses. These names are preserved as supplied references rather than converted into invented bibliographic entries.
\end{itemize}
\end{document}
